% Minimal retained preamble for lane-p2 cards taken from
% sources/logic-fl05-ramsey-like-theorems-for-the-schre/pinned-original-source.tex
% (arXiv:2412.11598v3, CC BY 4.0).  This source ships no separate preamble file:
% the pinned tex is a concatenation of the authors' main.tex and macros.tex and
% carries its preamble inline, so the blocks below are sliced mechanically from
% that file and nothing has been renamed or redefined.
%
% Copied verbatim from the source:
%   line 2        : the authors' class line
%   lines 3-5,7,9 : graphicx, amsthm, the ams bundle, cleveref, url
%   line 10       : tikz-cd, used by the authors' commutative diagrams
%   lines 20-108  : the authors' macro definitions, taken in full so that every
%                   macro of the retained spans is available
%   lines 113-124 : the authors' theorem declarations
%   lines 131-141 : the authors' machinery for restating a theorem
%   lines 143-148 : the authors' coloured comment macros
%
% Omitted source preamble lines and why:
%   line 1, 110-111 and 150 (concatenation comment, title, date, author): the
%       assembler supplies the title block, author list and date;
%   line 6 (hyperref) and line 8 (inputenc with the utf8x option): the forced
%       xelatex engine supplies the hyperref wiring itself and utf8x is
%       incompatible with it;
%   line 115 (numberwithin theorem by section): the assembler numbers the
%       retained material itself;
%   lines 120 and 121-123 are declarations that the authors left commented out or
%       that duplicate earlier ones and are not reproduced.
\documentclass{article}
\usepackage{graphicx}
\usepackage{amsthm}
\usepackage{amsfonts, amsmath, amssymb, amsthm}
\usepackage{cleveref}
\usepackage{url}
\usepackage{tikz-cd}
\newcommand{\Psf}{\mathsf{P}}
\newcommand{\Qsf}{\mathsf{Q}}
\newcommand{\Rsf}{\mathsf{R}}
\def\F{\mathcal{F}}
\def\U{\mathcal{U}}
\def\C{\mathcal{C}}
\def\I{\mathcal{I}}
\def\D{\mathcal{D}}
\def\L{\mathcal{L}}
\def\V{\mathcal{V}}
\def\A{\mathcal{A}}
\def\B{\mathcal{B}}
\def\P{\mathcal{P}}
\def\M{\mathcal{M}}
\def\Nc{\mathcal{N}}
\def\Q{\mathcal{Q}}
\newcommand{\cs}{2^{\NN}}
\newcommand{\bs}{\NN^{\NN}}
\newcommand{\bstr}{2^{<\NN}}
\newcommand{\uh}{\upharpoonright}
\newcommand{\converge}{\!\!\downarrow}
\newcommand{\diverge}{\!\!\uparrow}
\newcommand{\qvdash}{\operatorname{{?}{\vdash}}}
\newcommand{\nqvdash}{\operatorname{{?}{\nvdash}}}
\newcommand{\ISig}{\mathsf{I}\Sigma^0}
\newcommand{\Robinson}{\mathsf{Q}}
\newcommand{\IPi}{\mathsf{I}\Pi^0}
\newcommand{\IDelta}{\mathsf{I}\Delta^0}
\newcommand{\Iopen}{\mathsf{I_{open}}}
\newcommand{\BSig}{\mathsf{B}\Sigma^0}
\newcommand{\BPi}{\mathsf{B}\Pi^0}
\newcommand{\BDelta}{\mathsf{B}\Delta^0}
\newcommand{\BCSig}{\mathsf{BC}\Sigma^0}
\newcommand{\BCPi}{\mathsf{BC}\Pi^0}
\newcommand{\BCDelta}{\mathsf{BC}\Delta^0}
\newcommand{\LSig}{\mathsf{L}\Sigma^0}
\newcommand{\LPi}{\mathsf{L}\Pi^0}
\newcommand{\CSig}{\mathsf{C}\Sigma}
\newcommand{\RCA}[0]{\mathsf{RCA}}
\newcommand{\WKL}[0]{\mathsf{WKL}}
\newcommand{\WWKL}[0]{\mathsf{WWKL}}
\newcommand{\ACA}[0]{\mathsf{ACA}}
\newcommand{\ATR}[0]{\mathsf{ATR}}
\newcommand{\PIOOCA}[0]{\Pi^1_1\mathsf{-CA}}
\newcommand{\CAC}[0]{\mathsf{CAC}}
\newcommand{\ADS}[0]{\mathsf{ADS}}
\newcommand{\RRT}[0]{\mathsf{RRT}}
\newcommand{\FS}[0]{\mathsf{FS}}
\newcommand{\PFS}[0]{\mathsf{PFS}}
\newcommand{\TS}[0]{\mathsf{TS}}
\newcommand{\GRRT}[0]{\mathsf{GRRT}}
\newcommand{\LRRT}[0]{\mathsf{LRRT}}
\newcommand{\EM}[0]{\mathsf{EM}}
\newcommand{\SEM}[0]{\mathsf{SEM}}
\newcommand{\HalfEM}[0]{\mathsf{Half}\mbox{-}\mathsf{EM}}
\newcommand{\RT}[0]{\mathsf{RT}}
\newcommand{\SRT}[0]{\mathsf{SRT}}
\newcommand{\NN}[0]{\mathbb{N}}
\newcommand{\nat}[0]{\mathbb{N}}
\newcommand{\card}{\operatorname{card}}
\newcommand{\finsub}{\subseteq_{\mathtt{fin}}}
\newcommand{\W}{\mathrm{W}}
\newcommand{\sW}{\mathrm{sW}}
\newcommand{\lex}{\mathrm{lex}}
\newcommand{\base}{\operatorname{base}}
\newcommand{\range}{\operatorname{ran}}
\def\qt#1{``#1''}%

\newtheorem*{statement}{Statement}
\newtheorem{theorem}{Theorem}
\numberwithin{theorem}{section}
\newtheorem{maintheorem}[theorem]{Main Theorem}
\newtheorem{lemma}[theorem]{Lemma}
\newtheorem{question}[theorem]{Question}
\newtheorem{proposition}[theorem]{Proposition}
%\newtheorem{statement}[theorem]{Statement}
\newtheorem{remark}[theorem]{Remark}
\newtheorem{definition}[theorem]{Definition}
\newtheorem{corollary}[theorem]{Corollary}
\newtheorem{example}[theorem]{Example}

% Allow to repeat theorems
\makeatletter
\newtheorem*{rep@theorem}{\rep@title}
\newcommand{\newreptheorem}[2]{%
\newenvironment{rep#1}[1]{%
 \def\rep@title{#2 \ref{##1}}%
 \begin{rep@theorem}}%
 {\end{rep@theorem}}}
\makeatother

\newreptheorem{theorem}{Theorem}

% Comments
\usepackage{xcolor}
\newcommand{\ludovic}[1]{\textcolor{blue}{#1}}
\newcommand{\quentin}[1]{\textcolor{green}{#1}}
\newcommand{\lorenzo}[1]{\textcolor{orange}{#1}}
\newcommand{\oriola}[1]{\textcolor{red}{#1}}
\title{Two computable arrays of Turing indexes: a colouring from the k-th jump whose thin sets compute the k-th jump}
\author{Lorenzo Carlucci, Oriola Gjetaj, Quentin Le Houérou, Ludovic Levy Patey}
\date{arxiv:2412.11598v3; CC BY 4.0}
\begin{document}
\maketitle
\noindent\textit{Editorial note.} The counted unit is the article's own Lemma lem:arrays at source lines 683-685 together with its proof at source lines 687-697, a straight adjacent pairing in the pinned file. The pinned file is a concatenation of the authors' main and macro files; it is the article's fixed version arXiv:2412.11598v3 under CC BY 4.0.\par\medskip
\noindent\textit{Editorial note.} This source ships no separate preamble file, so the retained preamble is sliced mechanically from the pinned file by the lane and is shared, byte for byte, with the other cards of this source; it carries the authors' full macro block so that no macro of a retained span is left undefined. The first card of this source needed one revision for exactly that reason.\par\medskip
\noindent\textit{Editorial note.} The retained statement and proof contain no reference to another labelled result of the article and no citation, so this card carries no navigation pointer and no bibliography entry. The two arrays produced by the lemma are used by later results of the article, and those results are counted by other lanes; none of them is reproduced here.\par\medskip
\noindent\textit{Editorial note.} The proof's final sentence is a uniformity claim about the computation it has just built, and it is part of the author's own argument rather than a pointer to another result, so it is retained verbatim.\par\medskip
\noindent\textit{Allowed standard prerequisites (editorial).}\begin{itemize}
\item Thin sets, uniform moduli of a jump and computation from a set. A set is thin for a colouring of finite tuples in a colour when no tuple of the set of the relevant length has that colour; a uniform modulus of a set is a function that dominates every function computed by that set; and a set computes another when the latter is decidable relative to the former, written with the usual Turing-reduction symbol. The author's own notation is retained inside the counted statement and proof.. The counted lemma is the engine of the section: it produces both the colouring lived on the k-th jump and the indexes that turn any infinite thin set back into the k-th jump, and the following results of the section are built on exactly those two arrays.
\end{itemize}\medskip
\begin{lemma}\label{lem:arrays}
    There exists two computable arrays $(e_{n,k})_{n,k \in \NN}$ and $(d_{n,k})_{n,k \in \NN}$ of Turing indexes such that for every $n, k \in \NN^+$ with $n \geq 2$, $\Phi_{e_{n,k}}^{\emptyset^{(k)}}$ is a coloring $f_n^k : [\NN]^{n} \to n$ such that for every infinite $f_n^k$-thin set $H$, $\Phi_{d_{n,k}}^{H} = \emptyset^{(k)}$.
\end{lemma}

\begin{proof}
Let $k\in\NN^+$. Let $g_k$ be a uniform modulus of the set $\emptyset^{(k)}$. Note that $\Delta^0_{k+1}$-indexes for each function $g_k$ can be found computably, uniformly in~$k$. 

For $n \geq 2$ let $f_n^k(x_0, \dots,x_{n-1}) = n - 1$ if $g_k(x_0) \leq x_1$ and if $g_k(x_0) > x_1$, let $f_n^k(x_0, \dots,x_{n-1}) = i$ for $i < n - 1$ the largest value such that $g_k(x_0) > x_{i+1}$.

Let $H$ be an infinite $f_n^k$-thin set for some color $c$. The color $c$ cannot be $n-1$ as for a given $x_0 \in H$ there exists $x_1 <\dots < x_{n-1} \in H \setminus \{0, \dots, g_k(x_0)\}$, hence $f^k_n(x_0, \dots, x_{n-1}) = n-1$.

We claim $H$ is thin for color $n-2$. Indeed, assume by contradiction that there exists some tuple $x_0 < \dots < x_{n-1} \in H$ such that $g_k(x_0) > x_{n-1}$ so that $f^k_n$ takes color $n-2$ on $\{x_0, \dots, x_{n-1}\}$. Since $c < n-2$ and $H$ is infinite, there exists some $y_{c+2}< \dots < y_{n-1} \in H \setminus \{0, \dots ,g_k(x_0)\}$, therefore $f^k_n(x_0, \dots, x_{c+1},y_{c+2}, \dots, y_{n-1}) = c$, contradicting our assumption that $H$ is $f_n^k$-thin for $c$.

Write $H = \{x_0 < x_1 < \dots \}$. For every $i \in \NN$, $g_k(x_i) \leq x_{i + n - 1}$, hence $H$ computes a function dominating $g_k$ and therefore computes $\emptyset^{(k)}$. This computation can be done uniformly in the set $H$, $k$ and $n$. 
\end{proof}

\end{document}
